\documentclass[a4paper,14pt]{extarticle}
\usepackage[margin=0.8in]{geometry}
\usepackage{enumitem} % For customizing lists
\usepackage{graphicx}
\usepackage{amsmath}
\usepackage{makecell} % For making multi-line cells in tables
\usepackage{hyperref} % For hyperlinks
\usepackage{listings} % For code listings
\usepackage{xcolor}
\usepackage{easytable}
\usepackage{doi} % For adding doi to references
\usepackage{fontspec}

\setmonofont{IBM Plex Mono}
\setmainfont{IBM Plex Serif}

\definecolor{codegray}{rgb}{0.5,0.5,0.5}
\definecolor{codepurple}{rgb}{0.58,0,0.82}
\definecolor{backcolour}{rgb}{0.95,0.95,0.92}

\lstset{
  language=C,
  basicstyle=\normalsize\fontfamily{cmtt}\selectfont,
  keywordstyle=\color{blue},
  commentstyle=\color{green},
  breaklines=true,
  frame=single,
  numbers=none, % Remove line numbers
  numberstyle=\tiny\color{gray}, % Style for line numbers (not used here)
}


\newcommand{\unnumberedsection}[1]{
  \section*{#1}
  \addcontentsline{toc}{section}{#1}  
}

\newcommand{\unnumberedsubsection}[1]{
  \subsection*{#1}
  \addcontentsline{toc}{subsection}{#1}
}

\newcommand{\imagewithcaption}[3]{
  \begin{figure}[h!]
    \centering
    \includegraphics[width=#1\textwidth]{#2}
    \caption{#3}
    \label{fig:#2}
  \end{figure}
}


\begin{document}

\unnumberedsection{Satellite Project: Cloud Masking}

\unnumberedsection{Team Members}

\begin{center}
    \begin{TAB}(v,1cm,1cm)[5pt]{|l|l|l|l|}{|c|c|c|c|c|}% (rows,min,max)[tabcolsep]{columns}{rows}
    \textbf{ID} & \textbf{Name} & \textbf{Sec} & \textbf{BN} \\
    9213035 & Waleed Ousama Khamees & 2 & 30 \\
    9210684 & Ali Mohamed Farid & 1 & 25 \\
    9213347 & Mohamed Maher Hasan Amin & 2 & 12 \\
    9211164 & Mostafa Elsayed & 2 & 20 \\
    \end{TAB}
\end{center}

\unnumberedsection{Project Pipeline}
\imagewithcaption{1}{architecture.png}{Project Architecture}

\newpage

\unnumberedsection{PreProcessing}
The preprocessing pipeline includes:
\begin{itemize}
    \item Loading and organizing satellite images and their corresponding masks from TIFF files
    \item Converting GDAL datasets to NumPy arrays using 
     \\ \texttt{convert\_to\_nparr} function
    \item Splitting the dataset into training (60\%), validation (20\%), and test (20\%) sets
    \item Applying data augmentation with random horizontal flips, vertical flips, and 45-degree rotations
    \item Using PyTorch's DataLoader for efficient batch processing
\end{itemize}

\unnumberedsection{Exploratory Data Analysis}

\begin{itemize}
    \item Dataset consists of 10,573 satellite images with corresponding
    \\ cloud masks
    \item Images are in TIFF format with 4 spectral bands 
    \\ (shape: 4x512x512)
    \item Masks are binary (cloud vs. no-cloud) with shape 1x512x512
    \item Data distribution:
    \begin{enumerate}
        \item Training: 6,343 images
        \item Validation: 2,115 images
        \item Test: 2,115 images
    \end{enumerate}
    \item Dataset was shuffled with a fixed random seed (42) 
    \\ for reproducibility
\end{itemize}

\newpage

\unnumberedsection{Model Architecture}

We implemented a modified U-Net architecture with
\\ SqueezeNet-inspired Fire modules:
\begin{itemize}
    \item Encoder-decoder structure with skip connections
    \item Initial layer replaced with FireModule (squeeze-expand architecture)
    \item 4 levels of downsampling and upsampling
    \item Batch normalization and dropout (p=0.5) for regularization
    \item Final sigmoid activation for binary classification
    \item Input channels: 4 (for the 4 spectral bands)
    \item Output channels: 1 (binary mask)
\end{itemize}

Key components:
\begin{itemize}
    \item \texttt{FireModule}: Squeeze (1x1 conv) followed by parallel 1x1 and 3x3 expansions
    \item \texttt{ConvBlock}: Two convolutional layers with optional batch norm and dropout
    \item Transposed convolutions for upsampling
\end{itemize}

\unnumberedsection{Training Configuration}

\begin{itemize}
    \item Batch size: 16
    \item Learning rate: 0.001
    \item Number of epochs: 20
    \item Loss function: Binary Cross Entropy Loss (BCELoss)
    \item Optimizer: SGD
    \item Dropout probability: 0.5
    \item Training hardware: NVIDIA T4 GPU (via Kaggle)
    \item Early stopping based on validation Dice coefficient
\end{itemize}

\unnumberedsection{Training Results}

The model showed consistent improvement during training:
\begin{itemize}
    \item Initial training loss: 0.6924
    \item Final training loss: 0.5858
    \item Best validation Dice coefficient: 0.8562 (epoch 19)
    \item Training time: Approximately 4 hours for 20 epochs
\end{itemize}

\unnumberedsection{Evaluation Metrics}

\begin{itemize}
    \item Test Dice Coefficient: 0.8473
    \item Model performance was consistent between validation 
    \\ and test sets
    \item The Dice coefficient improved from 0.8027 in epoch 1 to 0.8562 in epoch 19
    \item Number of operations: 109.8 billion
    \item Model size: 7.7 million parameters
\end{itemize}

\unnumberedsection{Inference Pipeline}

\begin{itemize}
    \item Model loading and evaluation on test set
    \item Thresholding predictions at 0.5 for binary masks
    \item Batch processing of test images
    \item Generation of submission CSV file
\end{itemize}

\unnumberedsection{Future Work}

Potential improvements:
\begin{itemize}
    \item Experiment with different loss functions (Dice loss, Focal loss)
    \item Try more advanced architectures (Attention U-Net, ResUNet)
\end{itemize}

\end{document}
